\section{Experiments}
\label{sec:experiments}

We compare the balanced router with a dense feed-forward layer of equal cost and with unconstrained top-$k$ routing.\footnote{All numbers in this seed paper are illustrative and were not measured.}
Both tasks are synthetic: copying a sequence with long gaps, and predicting the parity of a sliding window.
Every model has four layers and model width $d = 256$; the mixture models use $n = 16$ experts with $k = 2$.

\begin{figure}[t]
  \centering
  \includegraphics[width=0.8\linewidth]{figures/routing.png}
  \caption{Expert load over training for unconstrained and balanced routing. The balanced router stays near the uniform load.}
  \label{fig:load}
\end{figure}

\Cref{fig:load} shows the load of the busiest expert during training.
Without a capacity limit, two experts receive most tokens within the first thousand steps; with $c = 1.25$ the busiest expert stays within the bound of \cref{prop:load}.
\Cref{tab:results} reports accuracy at the end of training.

\begin{table}[t]
  \centering
  \caption{Accuracy (\%) on the two synthetic tasks, mean of three runs (illustrative).}
  \label{tab:results}
  \begin{tabular}{@{}llcc@{}}
    \toprule
    \multicolumn{2}{@{}l}{Model} & Copy & Parity \\
    \midrule
    Dense & $d_\mathrm{ff} = 1024$ & 91.2 & 74.5 \\
    \multirow{2}{*}{Top-$k$} & unconstrained & 93.8 & 77.1 \\
     & balanced, $c = 1.25$ & 95.4 & 79.6 \\
    \bottomrule
  \end{tabular}
\end{table}

The balanced router improves on both baselines, and the gain is largest on the parity task, where unconstrained routing leaves most experts idle.
